\documentclass[10pt,a4paper]{amsart}
\usepackage[margin=19mm]{geometry}
\usepackage{amsmath,amssymb,mathtools}
\usepackage[scr=rsfs,cal=euler]{mathalfa}
\usepackage{xfrac}
\usepackage[unicode,hidelinks,bookmarks=false]{hyperref}
\newcommand{\ie}{i.e.\ }
\newcommand{\cf}{cf.\ }
\newcommand{\resp}{respectively\ }
\newcommand{\Subsection}{\subsection}
\providecommand{\nobreakdash}{\nobreak\hbox{-}}
\setlength{\emergencystretch}{2em}
\multlinegap0pt
\usepackage[shortlabels]{enumitem}
\usepackage{xfrac}
\usepackage{bm}
\usepackage{nicefrac}
\usepackage{amssymb}
\usepackage{dsfont}
\usepackage{multirow}
\usepackage{xcolor}
\usepackage{booktabs}
\usepackage{mathtools}
\usepackage{caption}
\newtheorem{assumption}{Assumption}
\newtheorem{proposition}{Proposition}
\newcommand{\N}{\mathbb{N}}
\newcommand{\R}{\mathbb{R}}
\newcommand{\bbP}{\mathbb{P}}
\newcommand{\bbR}{\mathbb{R}}
\newcommand{\cF}{\mathcal{F}}
\newcommand{\from}{\colon}
\newcommand{\id}{\operatorname{Id}}
\newcommand{\norm}[2]{     \| #1       \|_{ #2 }}
\newcommand{\rd}{\mathop{}\!\mathrm{d}}
\title{Multiple and weak Markov properties in Hilbert spaces with applications to fractional stochastic evolution equations}
\author{Kristin Kirchner, Joshua Willems}
\date{Source: arXiv:2310.13536v1 (CC BY 4.0)}
\begin{document}
\maketitle

\noindent\emph{Source and licence.} Multiple and weak Markov properties in Hilbert spaces with applications to fractional stochastic evolution equations. Version arXiv:2310.13536v1 (CC BY 4.0), distributed by arXiv under the Creative Commons Attribution 4.0 International licence, http://creativecommons.org/licenses/by/4.0/, as recorded in the arXiv licence metadata for this exact version and shown on the abstract page https://arxiv.org/abs/2310.13536v1. Full licence notice: \texttt{licenses/arxiv-2310.13536-CC-BY-4.0.txt}.\par\smallskip

\noindent\textit{Editorial scope.} Counted unit: the article's proposition prop:recurrence-zeta (source0.tex lines 2470--2492). The article's complete original proof of it is delivered with it (source0.tex lines 2603--2685, one \texttt{proof} environment of 83 lines, which the article places away from the statement). No mathematics is written, repaired or re-derived: the statement and the proof are the article's own lines, in the article's own notation, and no retained line is substituted or re-flowed.  Retained uncounted as the source-local context the proof needs: lines 1797--1809; lines 1814--1828; lines 1852--1855; lines 2363--2372; lines 2399--2404.  Nothing was omitted from any retained span, so the package carries no omission record.  No retained line contains a citation, so the wrapper carries no bibliography and no bibliography entry.  No retained line contains a figure environment, an image-inclusion command or drawing code, so no image is delivered and the package declares no asset dependency.  Editorial formatting only, in addition: an AMS \texttt{amsart} preamble that restates the article's own declarations and macros used by the retained lines, namely the environments \texttt{assumption}, \texttt{proposition} on separate counters, together with the standard packages those lines need.  The retained environments print in this document as Assumption 1, Proposition 1 in order of appearance, whereas the article prints the same items with the numbering of its own full document.  The complete native source of the article as distributed in the arXiv e-print for this exact version is bundled unchanged as \texttt{sources/148805/source0.tex}.
\begin{assumption}\label{ass:standing}
	The linear operator
	$-A \from \mathsf D(A) \subseteq U \to U$
	on the separable real Hilbert space $U$ generates an
	\emph{exponentially stable} $C_0$-semigroup $(S(t))_{t\ge0}$ of bounded
	linear operators on $U$, i.e.,
	%
	\begin{equation}\label{eq:semigroup-est}
		\exists M_0 \in [1,\infty), w \in (0, \infty) : \forall t \in [0,\infty) : \quad
		\norm{S(t)}{\mathscr L(U)} \le M_0 e^{-wt}.
	\end{equation}
	%
\end{assumption}

\begin{assumption}\label{ass:HAQ}
	$\,$
	%
	\begin{enumerate}[leftmargin=1cm]
		\item \label{ass:HAQ:1}
		There exist $\gamma_0 \in (\nicefrac{1}{2}, \infty)$ and $Q \in \mathscr L^+(U)$ such that
		%
		\[
		\int_0^\infty \| t^{\gamma_0-1} S(t) Q^\frac{1}{2} \|_{\mathscr L_2(U)}^2 \rd t < \infty.
		\]
		\item \label{ass:HAQ:2}
		The $C_0$-semigroup
		$(S(t))_{t\ge0}$ is analytic.
	\end{enumerate}
\end{assumption}

\begin{equation}\label{eq:analytic-semigroup-est}
	\exists M_j \in [1,\infty) \colon \: \forall t \in (0,\infty) : \quad
	\norm{A^j S(t)}{\mathscr L(U)} \le M_j t^{-j} e^{-wt}.
\end{equation}

\begin{equation}\label{eq:def-incgamma}
	\begin{aligned}
		\overline\Gamma(n, tA) 
		:= 
		\begin{cases}
			\sum_{j=0}^{n-1} \frac{t^j}{j!} A^j S(t), \quad &t \in (0,\infty); \\
			\id_U, \quad &t = 0,
		\end{cases}
	\end{aligned}
\end{equation}

	\begin{equation}\label{eq:def-zetaN}
		\zeta_{N}(t \mid t_0)\bm\xi 
		:= 
		\sum_{k=0}^{N-1} \frac{(t-t_0)^k}{k!}
		\overline\Gamma(N-k, (t-t_0)A)\, \xi_k,
	\end{equation}

\begin{proposition}\label{prop:recurrence-zeta}
	Let $N \in \{2,3,\dots\}$, $t_0 \in \bbR$ and $\bm\xi \in L^2(\Omega, \cF_{t_0},\bbP; U^{N})$
	be given,
	where $\xi_k \in L^2(\Omega; \mathsf D(A))$ for $k \in \{0, \dots,  N-2\}$.
	Suppose that Assumptions~\textup{\ref{ass:standing} and~\ref{ass:HAQ}\ref{ass:HAQ:2}} hold.
	Then the process $(\zeta_{N}(t \mid t_0)\bm{\xi})_{t \in [t_0,\infty)}$ from~\eqref{eq:def-zetaN}
	is infinitely mean-square differentiable at any $t \in (t_0, \infty)$
	and, for $n \in \{0,\dots,N-2\}$,
	%
	\begin{equation}\label{eq:recurrence-zeta}
		\left(\tfrac{\mathrm d}{\mathrm dt} + A\right)\tfrac{\mathrm d^n}{\mathrm dt^n}
		\zeta_{N}(t \mid t_0) \bm \xi
		= 
		\tfrac{\mathrm d^n}{\mathrm dt^n} \zeta_{N-1}(t \bigm\vert t_0)(\xi_{k+1} + A\xi_{k})_{k=0}^{N-2},
		\quad
		\bbP\text{-a.s.}
	\end{equation}
	%
	Moreover, we have $\left(\tfrac{\mathrm d}{\mathrm dt} + A\right) \tfrac{\mathrm d^n}{\mathrm dt^n} \zeta_{1}(t \mid t_0)\xi 
	= 
	0$, $\bbP$-a.s.,
	for $\xi \in L^2(\Omega, \cF_{t_0}, \bbP; U)$.
\end{proposition}

\begin{proof}[Proof of Proposition~\ref{prop:recurrence-zeta}]
	We first make some general remarks regarding the operators $\overline\Gamma(n, tA)$
	from~\eqref{eq:def-incgamma}.
	Under Assumption~\ref{ass:HAQ}\ref{ass:HAQ:2},
	estimate~\eqref{eq:analytic-semigroup-est} implies that
	the set $\{t^j A^jS(t) : t \in (0,\infty)\} \subseteq \mathscr{L}(U)$ is uniformly bounded.
	It follows that
	$t \mapsto \overline\Gamma(n, tA)$ 
	is a strongly continuous
	function from $[0,\infty)$ to $\mathscr{L}(U)$ for any $n \in \N$,
	which at $t \in (0,\infty)$ admits a classical derivative
	satisfying the recurrence relation
	%
	\begin{equation}\label{eq:recurrence-incompletegamma}
		\left(\tfrac{\mathrm d}{\mathrm d t} + A\right)\overline{\Gamma}(n, tA) 
		= 
		\begin{cases}
			0, &n = 1; \\
			A\overline{\Gamma}(n-1, tA), &n \in \{2,3,\dots\}.
		\end{cases}
	\end{equation} 
	%
	To prove the proposition, we may assume $t_0 = 0$, so fix $t \in (0,\infty)$.
	For arbitrary $M \in \N$, $j \in \N_0$ and $\bm\eta \in L^2(\Omega; U^{M})$, we define
	%
	$
	\zeta_{M,j}(t) \bm\eta
	:= 
	A^j \zeta_{M}(t \mid 0) \bm\eta.
	$
	%
	Combining the product rule in the form 
	$(\frac{\mathrm d}{\mathrm dt} + A)(uv) = u'v + u[(\frac{\mathrm d}{\mathrm dt} + A)v]$
	with the above recurrence relation yields for $M \in \N$
	%
	\begin{align}
		(\tfrac{\mathrm d}{\mathrm dt} &+ A)\zeta_{M,j}(t) \bm \eta
		\\&=
		\sum_{k=1}^{M-1} \frac{t^{k-1}}{(k-1)!} A^j \overline\Gamma(M-k, tA) \eta_k
		+
		\sum_{k=0}^{M-2} \frac{t^k}{k!} A^{j+1} \overline\Gamma(M-1-k, tA) \eta_k
		\notag
		\\
		&=
		\zeta_{(M-1),j}(t) (\eta_{k+1})_{k=0}^{M-1} + \zeta_{(M-1),(j+1)}(t) (\eta_{k})_{k=0}^{M-1}
		\label{eq:recurrence-Z0-SP-k0j}.
	\end{align}
	%
	This shows in particular that~\eqref{eq:recurrence-zeta}
	holds for integers $N \ge 2$ and $n = 0$, 
	by applying~\eqref{eq:recurrence-Z0-SP-k0j} with $M = N$, $\bm\eta = \bm\xi$ and $j = 0$.
	Moreover, note that
	%
	\[
	\left(\tfrac{\mathrm d}{\mathrm dt} + A\right) \zeta_{1,j}(t) \, \eta 
	=
	\left(\tfrac{\mathrm d}{\mathrm dt} + A\right) A^j S(t) \, \eta
	=
	0.
	\]
	%
	Iteratively applying~\eqref{eq:recurrence-Z0-SP-k0j} and the latter identity then
	yields that $\zeta_{M,j}\bm\eta$ is $M-1$ times (mean-square) differentiable
	with an $n$th derivative of the form
	%
	\begin{equation}\label{eq:deriv-initial-val-term}
		\zeta^{(n)}_{M,j}(t) \bm \eta 
		=
		\sum_{\ell=0}^n \sum_{m=0}^\ell C_{\ell,m} \zeta_{(M-\ell),(j+n-m)}(t) B_{\ell, m} \bm \eta,
	\end{equation}
	%
	where $C_{\ell,m} \in \R$, $B_{\ell, m} \in \mathscr L(U^{M}; U^{M-\ell})$ and
	$\zeta_{(M-\ell),(j+n-m)} := 0$ if $M-\ell < 1$.
	In particular, $\zeta_{N}(\,\cdot \mid t_0)\bm{\xi}$ is $N-1$ times (mean-square) differentiable as claimed.
	In order to deduce that~\eqref{eq:recurrence-zeta} also holds for $n \in \{1, \dots, N-2\}$, we need
	to justify taking the $n$th derivative on both sides and commuting
	it with $A$. Since $A$ is closed,
	it suffices to verify that $\zeta_{N}'(\,\cdot \mid t_0)\bm \xi$,
	$A^j \zeta_{N}(\,\cdot \mid t_0)\bm \xi$,
	$\zeta_{N-1}(\,\cdot \mid t_0)  (\xi_{j+1})_{j=0}^{N-1}$ and $A^j \zeta_{N-1}(\,\cdot \mid t_0)  (\xi_{j})_{j=0}^{N-1}$
	admit $n$th derivatives for $j \in \{0,1\}$. Indeed,
	these assertions follow from~\eqref{eq:deriv-initial-val-term}.
\end{proof}

\end{document}
